\documentclass[UTF8]{ctexart}

%\RequirePackage{amsmath} \RequirePackage{amssymb}
\usepackage{amscd,amsthm,amsfonts,amssymb,amsmath}
\usepackage[colorlinks=true]{hyperref}
%\usepackage{fullpage}
\usepackage[all]{xy}
\usepackage{mathrsfs}
\usepackage{tikz,mathpazo}
%\usetikzlibrary{quotes,angles}
%\usetikzlibrary{calc}
%\usetikzlibrary{decorations.pathreplacing}
\usepackage{bm}
\usepackage{geometry}
\usepackage{xcolor}
\usepackage{eso-pic}
\usepackage{CJK}
\usepackage{arydshln}
\usepackage{mathptmx}
\usepackage[11pt]{moresize}



\newcommand{\BC}{{\mathbb {C}}}

\newcommand{\BN}{{\mathbb {N}}}
\newcommand{\BQ}{{\mathbb {Q}}}
\newcommand{\BR}{{\mathbb {R}}}
\newcommand{\BZ}{{\mathbb {Z}}}
\newcommand{\bv}{{\bm v}}
\newcommand{\bw}{{\bm w}}
\newcommand{\bu}{{\bm u}}
\newcommand{\bx}{{\bm x}}
\newcommand{\by}{{\bm y}}

\newcommand{\bb}{{\bm b}}
\newcommand{\Coker}{{\mathrm{Coker}}}


\newcommand{\End}{{\mathrm{End}}}
\newcommand{\GL}{{\mathrm{GL}}}

\newcommand{\SO}{{\mathrm{GSO}}}
\newcommand{\Hom}{{\mathrm{Hom}}}

\renewcommand{\Im}{{\mathrm{Im}}}

\newcommand{\Isom}{{\mathrm{Isom}}}
\newcommand{\Id}{{\mathrm{Id}}}



\newcommand{\Ker}{{\mathrm{Ker}}}



\newcommand{\rank}{{\mathrm{rank}}}


\newcommand{\SU}{{\mathrm{SU}}}
\newcommand{\Sym}{{\mathrm{Sym}}}

\newcommand{\sub}{{\mathrm{sub}}}

\newcommand{\tr}{{\mathrm{tr}}}

\newcommand{\matrixx}[4]{\begin{pmatrix}
		#1 & #2 \\ #3 & #4
	\end{pmatrix} }
	
	
	
	\newcommand{\wt}{\widetilde}
	\newcommand{\wh}{\widehat}
	\newcommand{\pp}{\frac{\partial\bar\partial}{\pi i}}
	\newcommand{\pair}[1]{\langle {#1} \rangle}
	\newcommand{\wpair}[1]{\left\{{#1}\right\}}
	\newcommand{\intn}[1]{\left( {#1} \right)}
	\newcommand{\norm}[1]{\|{#1}\|}
	\newcommand{\sfrac}[2]{\left( \frac {#1}{#2}\right)}
	\newcommand{\ds}{\displaystyle}
	\newcommand{\ov}{\overline}
	\newcommand{\incl}{\hookrightarrow}
	\newcommand{\lra}{\longrightarrow}
	\newcommand{\lla}{\longleftarrow}
	\newcommand{\ra}{\rightarrow}
	\newcommand{\imp}{\Longrightarrow}
	\newcommand{\lto}{\longmapsto}
	\newcommand{\bs}{\backslash}
	
	
	
	
	\newtheorem{thm}{定理}
	\newtheorem{cor}[thm]{推论}
	\newtheorem{lem}[thm]{引理}
	\newtheorem{prop}[thm]{命题}
	\newtheorem{defn}[thm]{定义}
	\newtheorem{ques}[thm]{问题}
	\newtheorem{eg}[thm]{例题}
	\newtheorem{ex}[thm]{习题}
	\newtheorem{rmk}[thm]{注释}
		\newtheorem{ntn}[thm]{记号}

	\newcommand{\sk}{\medskip}\newcommand{\bsk}{\bigskip}
	\newcommand{\s}{\sk\noindent}\newcommand{\bigs}{\bsk\noindent}
	
	
	%accented words

	\def\mat(#1,#2,#3,#4){
		\left[\begin{array}{cc}
			#1 & #2 \\ #3 & #4\end{array}
		\right]
	}
		\def\clm(#1,#2){
		\left[\begin{array}{c}
			#1 \\ #2 \end{array}
\right]
	}
		\def\clmn(#1,#2){
		\left[\begin{array}{c}
			#1 \\ \vdots\\ #2 \end{array}
\right]
	}
		\def\clmt(#1,#2,#3){
		\left[\begin{array}{c}
			#1 \\ #2  \\ #3\end{array}
\right]
	}
		\def\clmf(#1,#2,#3){
		\left[\begin{array}{c}
			#1 \\ #2  \\ \vdots \\ #3\end{array}
\right]
	}


		
\begin{document}


\title{习题课材料（七）}

\date{}

\maketitle

\vspace{-1cm}



\noindent{\Large\textbf{注: 带 $\heartsuit $号的习题有一定的难度、比较耗时, 请量力为之.}}


\begin{ex}
若$A_1, A_2, \ldots, A_{2020}$是$2020$个$2019$阶方阵，求证: 则关于$\mathbf{x}=(x_1,\ldots ,x_{2020})$的方程$$\det \left(x_1A_1+x_2A_2+\cdots +x_{2020}A_{2020}\right)=0$$存在非零解。
\end{ex}

\begin{ex}[$\heartsuit $]
	计算行列式
$$\det A=\left|\begin{array}{ccccccc}
                 1 & 2 & 3 & \cdots & n-2 & n-1 & n \\
                 2 & 3 & 4 & \cdots & n-1 & n & n \\
                 3 & 4 & 5 & \cdots & n & n & n \\
                 \vdots & \vdots & \vdots &  & \vdots & \vdots & \vdots \\
                 n-1 & n & n & n & n & n & n \\
                 n & n & n & n & n & n & n
               \end{array}
\right|$$
\end{ex}



\begin{ex}[$\heartsuit $]
计算行列式
$$\det A=\left|\begin{array}{cccccc}
                 1 & 2 & 3 & \cdots & n-1 & n \\
                 n & 1 & 2 & \cdots & n-2 & n-1 \\
                 n-1 & n & 1 & \cdots & n-3 & n-2 \\
                 \vdots & \vdots & \vdots &  & \vdots & \vdots \\
                 2 & 3 & 4 & \cdots & n & 1
               \end{array}
\right|$$
\end{ex}


\begin{ex}
	设 $n$阶方阵 $A$的第 $i$行第 $j$列的元素为 $\min(i,j)$. 求 $\det A$.
\end{ex}

[记号] 设$A=(a_{ij})$是$n$阶方阵,（$n\geq 2$）,$C=(C_{ij})$，其中$C_{ij}$为$a_{ij}$的代数余子式
\begin{ex}
设
\begin{eqnarray*}
D=\left |
\begin{array}{cccc}
 1 &-1 &1 &2\\
 2 &1 &0 &-1\\
 -2 & 2 &-2 &-3\\
 -1 &2 &-2 &-3
\end{array} \right|
\end{eqnarray*}
不直接计算$C_{ij}$, 求解以下各题：

(1) $-2C_{11}+2C_{21}+3C_{31}+4C_{41};$

(2) $C_{13}+C_{23}+C_{33}+C_{43}.$
\end{ex}

\begin{ex}
设
\begin{eqnarray*}
 D=\left |
\begin{array}{cccc}
1 &-1 &1 &2 \\ 3 &6 &1 &1\\ 1 &2 &2 &2\\ -1 &2 & -2& -3
\end{array}
\right |
\end{eqnarray*} 求$S_1=C_{21}+2C_{22}$ 和$S_2=C_{23}+C_{24}$.
\end{ex}



\begin{ex}
设
\begin{eqnarray*}
D(x)=\left|
\begin{array}{ccc}
a_{11}(x) &a_{12}(x) &a_{13}(x)\\
a_{21}(x) &a_{22}(x) &a_{23}(x)\\
a_{31}(x) &a_{32}(x) &a_{33}(x)
\end{array}
\right|.
\end{eqnarray*}
其中函数$a_{ij}$可导.

(1)求证:\begin{eqnarray*}
D'(x)&=\left|
\begin{array}{ccc}
a'_{11}(x) &a'_{12}(x) &a'_{13}(x)\\
a_{21}(x) &a_{22}(x) &a_{23}(x)\\
a_{31}(x) &a_{32}(x) &a_{33}(x)
\end{array}
\right|+\left|
\begin{array}{ccc}
a_{11}(x) &a_{12}(x) &a_{13}(x)\\
a'_{21}(x) &a'_{22}(x) &a'_{23}(x)\\
a_{31}(x) &a_{32}(x) &a_{33}(x)
\end{array}
\right|\\&+\left|
\begin{array}{ccc}
a_{11}(x) &a_{12}(x) &a_{13}(x)\\
a_{21}(x) &a_{22}(x) &a_{23}(x)\\
a'_{31}(x) &a'_{32}(x) &a'_{33}(x)
\end{array}
\right|.
\end{eqnarray*}

(2)将上述结论推广到$n$阶的情形

\end{ex}


\begin{ex}
	设$A$为可逆方阵， $D$为方阵， 证明：
\begin{eqnarray*}
\left|
\begin{array}{cc}
A &B\\
C &D\\
\end{array}
\right|=|A||D-CA^{-1}B|.
\end{eqnarray*}
	\end{ex}



\begin{ex}[$\heartsuit $]
求如下的$n$阶范德蒙行列式：
\begin{eqnarray*}
\left|
\begin{array}{ccccc}
1 & x_1 & \cdots & x_1^{n-1}  \\
1 & x_2 & \cdots & x_2^{n-1}\\
\vdots & \vdots & \vdots & \vdots \\
1 & x_n & \cdots & x_n^{n-1}
\end{array}
\right|.
\end{eqnarray*}
\end{ex}
\begin{proof}[提示]参考《线性代数入门》p.164-165
\end{proof}

\begin{ex}[$\heartsuit $]
设$Q$是$n$阶正交矩阵，即$Q^TQ=QQ^T=I_n$。

(1) 若$|Q|<0$，求证：$|Q+I_n|=0$，因此存在非零向量$\mathbf{v} \in \mathbb{R}^n$，使得$Q\mathbf{v}=-\mathbf{v}$。

(2) 设 $|Q|>0$. 证明当$n$是奇数时等式$|Q-I_n|=0$总成立. 当n是偶数时, 判断等式$|Q-I_n|=0$是否成立. 若成立, 请给出证明;若不成立, 请举出反例.
\end{ex}



\end{document}
